%------------------------------------------------------------------------------
% \PART 5
%------------------------------------------------------------------------------
\chapter[Аппаратная реализация IP-блока нейронной сети для распознавания рукописных цифр]
{АППАРАТНАЯ РЕАЛИЗАЦИЯ IP-БЛОКА НЕЙРОННОЙ СЕТИ ДЛЯ РАСПОЗНАВАНИЯ РУКОПИСНЫХ ЦИФР}

%------------------------------------------------------------------------------
% \PART 5.1 Text
%------------------------------------------------------------------------------
\section{Описание пакета Xilinx Vivado}
\hspace*{12.5 mm}Xilinx Vivado – это современный программный пакет для 
проектирования цифровых схем на базе программируемых логических интегральных 
схем, разработанный компанией Xilinx. Данный программный комплекс 
является преемником Xilinx ISE и значительно превосходит его по возможностям и 
производительности. 

Vivado предназначен для автоматизированного проектирования 
цифровых схем, начиная с этапа описания логики на языках описания аппаратуры 
и заканчивая загрузкой сгенерированного битового потока в 
устройство\cite{XilinxUG}.  

Одним из ключевых преимуществ Vivado является использование более 
производительных алгоритмов синтеза и оптимизации, что позволяет существенно 
ускорить процесс разработки. В отличие от традиционных средств проектирования, 
в Vivado реализована поддержка параллельных вычислений, что делает его особенно
эффективным при работе с большими проектами.  

В состав пакета Vivado входят несколько инструментов, среди 
которых\cite{XilinxUG}:  

  1 {Vivado IDE} – интегрированная среда разработки, предоставляющая удобный 
графический интерфейс для управления проектами, анализа схемотехнической логики
и моделирования.

  2 {Vivado Synthesis} – инструмент синтеза логики, который преобразует код HDL
в оптимизированное представление на уровне вентилей.

  3 {Vivado Implementation} – модуль размещения, трассировки и оптимизации 
логики на кристалле FPGA.\@

  4 {Vivado Simulator} – встроенный инструмент для   моделирования цифровых 
схем на уровне поведенческой, функциональной и временной верификации.

  5 {Vivado IP Integrator} – инструмент для работы с IP-ядрами и объединения их
в сложные системы на кристалле.

  6 {Vivado High-Level Synthesis (HLS)} – средство высокоуровневого синтеза, 
позволяющее описывать аппаратные алгоритмы на языках C/C++ и автоматически 
преобразовывать их в оптимизированный RTL-код.

Основные возможности Vivado:  

  1 Поддержка языков описания аппаратуры VHDL и Verilog. Vivado позволяет 
разрабатывать схемы на двух основных языках HDL, а также поддерживает 
SystemVerilog для верификации цифровых проектов. Это дает широкие 
возможности для разработки и тестирования аппаратуры.

   2 Автоматизированный синтез и размещение логики на FPGA.\@ Vivado включает 
мощный механизм синтеза, который автоматически преобразует HDL-код в 
оптимизированную сетевую схему, выполняет ее размещение на кристалле FPGA и 
прокладывает соединения между логическими блоками.

  3 Встроенные инструменты для моделирования и отладки. 

  4 Генерация IP-ядер и использование встроенных библиотек компонентов. 


Для реализации нейронной сети Vivado используется на нескольких ключевых этапах
разработки. Во-первых, он применяется для генерации аппаратной логики, 
необходимой для выполнения расчетов внутри FPGA.\@ Во-вторых, с его помощью 
производится компиляция проекта и синтез аппаратной схемы. В-третьих, Vivado 
используется для загрузки сгенерированного битового потока в ПЛИС и отладки 
работы модели.  

%------------------------------------------------------------------------------
% \PART 5.2 Text
%------------------------------------------------------------------------------
\section{Описание функциональных блоков}
\hspace*{12.5 mm}Аппаратная реализация нейронной сети включает несколько 
ключевых функциональных блоков:

  1 Блок обработки пикселя.

  2 Блок вычисления тангенса.

  3 Блок вычисления softmax.

  4 Блок памяти изображения.

  5 Блок управляющего автомата.

  6 Блок MAC.

Каждый из этих блоков реализуется в виде аппаратного модуля на HDL.\@

В блоке обработки пикселя осуществляется установка параметра смещения для 
подготовки к вычислению слоя, а также выбирается соответствующий данному этапу 
весовой коэффициент из памяти.

В блоке вычисления тангенса осуществляется аппроксимированный расчет функции 
активации тангенс гиперболический для LST-1 слоя.

Прямая реализация вычисления тангенса будет занимать большое количество 
вычислительных ресурсов. Поэтому для удобства реализации используется 
аппроксимация в соответствии с формулой\cite{TANH}:
\vspace{1mm}
\begin{equation}
    \tanh{(x)}\approx F(x) = 
      \begin{cases}
      \mathrm{sign}(x),         & \mathrm{\text{если}}\;\;  |x|>2, \\
      (1 + \frac{x}{4})\cdot x, & \mathrm{\text{если}}\;\; -2<x<0, \\
      (1 - \frac{x}{4})\cdot x, & \mathrm{\text{если}}\;\;  0<x<2. \\
      \end{cases}
\end{equation}
\vspace{1mm}

В блоке вычисления softmax происходит сравнение входных данных и определение 
наиболее вероятного класса.

В блоке памяти изображения храниться результат обработки на каждом этапе. От 
приемки входных данных, до выхода LST-1 слоя.

Блок управляющего автомата осуществляет управление и синхронизацию между всеми 
блоками устройства и обрабатывает входные управляющие сигналы.

Блок MAC строится на основе блока DSP48. Данный блок используется для 
вычисления операции умножения с накоплением. В данном блоке по сигналу 
инициализации устанавливается смещение, необходимо для текущего этапа 
вычислений, а затем выполнение операции MAC на каждом такте синхронизации.

Реализация данных блоков на языке SystemVerilog представлена в приложении Г.

%------------------------------------------------------------------------------
% \PART 5.3 Text
%------------------------------------------------------------------------------
\section{Описание интерфейса нейронной сети}

\hspace*{12.5 mm}Как было показано ранее на рисунке 4.5 нейронная сеть состоит
из пяти обязательных входных и двух выходных сигналов. 
Взаимодействие процессорной системы и программируемой логики 
осуществляется через интерфейс AXI4-lite. 

Необходимым условием, выдвигаемым к IP-блоку будет поддержка 
интерфейса AXI4-lite. Для упрощения формирования IP-блока 
будем использовать связку AXI-uP. uP интерфейс напрямую 
взаимодействует с модулем верхнего уровня. Далее данные 
передаются в AXI-slave модуль, который и подключается к 
процессорной системе.

%------------------------------------------------------------------------------
% \PART 5.4 Text
%------------------------------------------------------------------------------
\section{Описание процесса создания блочной диаграммы}
\hspace*{12.5 mm}Рассмотрим основные задачи связанные с созданием блочной 
диаграммы (Block diagram), подключением готовых IP-ядер и созданием собственных, 
а также созданием топ-модуля на языке Verilog. Создадим блочную диаграмму, что 
показано на рисунке~\ref{fig:creat-bd}.

\insertfigure{creat-bd}{pic/create_bd.png}{Создание блочной диаграммы}{9cm}

Приступим к добавлению выводов процессорной системы в ПЛИС и сопутствующих 
IP-ядер, для этого нужно добавить ZYNQ7. После САПР Vivado предложит запустить 
мастер автоматического подключения <<Block Automation>>, что показано на 
рисунке~\ref{fig:sel-bd}.

\insertfigure{sel-bd}{pic/sel_ip.png}{Окно выбора IP-ядра и блочной диаграммы создаваемого проекта с <<Мастером подключений>> типа <<Block Automation>>}{16cm}

Мастер автоматического подключения <<Block Automation>> ZYNQ7 обладает опциями,
показанными на рисунке~\ref{fig:master-bd}.

\insertfigure{master-bd}{pic/mast-set.png}{Опции мастера автоматической настройки ZYNQ7}{15cm}

Перейдем к созданию собственого IP-ядра, что показано на 
рисунке~\ref{fig:ip-create}

\insertfigure{ip-create}{pic/create-ip.png}{Создание IP-ядра в САПР Vivado}{15cm}

Создадим новое IP-ядро и выберем тип интерфейса AXI4(рисунок~\ref{fig:ip-axi}).

\insertfigure{ip-axi}{pic/create-ip-axi.png}{Выбор типа IP-ядра}{14.5cm}

Зададим имя и описание нового IP-ядра(рисунок~\ref{fig:ip-name}).

\insertfigure{ip-name}{pic/ip-name.jpg}{Описание IP-ядра}{14.5cm}

Укажем типы интерфейсов будущего IP-ядра, а также задаим имя AXI4-Lite 
интерфейса как s\_axi, что показано на рисунке~\ref{fig:s_axi}

\insertfigure{s_axi}{pic/s-axi.jpg}{Аппаратные интерфейсы IP-ядра}{14cm}

Финализируем создание IP-ядра и выберем его редактирование(рисунок~\ref{fig:ip-final}).

\insertfigure{ip-final}{pic/ip-final.jpg}{Завершение создания IP-ядра}{14cm}


Далее выполним добавление всех файлов в IP-блок, что 
представлено на рисунке~\ref{fig:add-files}

\insertfigure{add-files}{pic/add_files.jpg}{Добавление файлов}{14cm}

Затем ядро можно перепаковать, что представлено на рисунке~\ref{fig:add-files}

\insertfigure{ip-last}{pic/ip_last.jpg}{Завершение создания IP-ядра}{14cm}

Далее необходимо перейти к построению блочной диаграммы аппаратной части 
проекта в среде Vivado. Блочная диаграмма представляет собой графическое 
описание аппаратных соединений между IP-ядрами, процессорной системой и 
периферией.

Для обеспечения связи между несколькими IP-блоками и процессором в систему 
добавляется модуль AXI Interconnect. Этот модуль играет ключевую роль в 
маршрутизации данных, поступающих от мастер-устройств к одному или нескольким 
ведомым (slave) IP-блокам. Он автоматически адаптирует адресацию, сигналы 
управления и синхронизации, обеспечивая корректный обмен данными в системе.

Также обязательным элементом системы является Processor System Reset. Этот блок 
управляет сигналами сброса всех модулей, входящих в состав блочной диаграммы. 
Он обеспечивает корректную и синхронную инициализацию всех компонентов после 
подачи питания или программного сброса, предотвращая возможные ошибки при 
старте системы.

На рисунке~\ref{fig:bd} показан результирующий вид блочной диаграммы, 
содержащей процессорную систему, пользовательский IP-блок, AXI Interconnect и 
блок сброса. Все соединения между модулями реализованы с соблюдением требований 
интерфейсов AXI и рекомендаций Vivado.

\insertfigure{bd}{pic/bd}{Результирующий вид блочной диаграммы}{16cm}

Для корректной работы системы требуется обеспечить стабильный тактовый сигнал с 
частотой 80 МГц. По умолчанию, тактовый генератор в Zynq-процессорной системе 
может выдавать несколько частот, но для нужд пользовательского IP-блока 
необходима именно частота 80 МГц.

Для этого на вкладке Clock Configuration в настройках Processing System 7 
выполняется установка пользовательской частоты. Один из выходных тактовых 
сигналов (FCLK\_CLK0) настраивается на значение 80 МГц. Этот сигнал затем 
подается на все соответствующие модули системы, включая пользовательский 
IP-блок, обеспечивая их синхронную работу.

Процесс настройки частоты подробно продемонстрирован на рисунке~\ref{fig:clk}.

\insertfigure{clk}{pic/clk}{Установка частоты 80 МГц в блоке Processing System}{16cm}


%------------------------------------------------------------------------------
% \PART 5.5 Text
%------------------------------------------------------------------------------
\section{Результаты синтеза и симуляции}
\hspace*{12.5 mm}Для проверки работоспособности написанного устройства был 
составлен тест, в котором формируются управляющие воздействия и изображение. 
Работа устройства начинается по сигналу start\_i. 

Для тестирования было взято два изображения, что показано на 
рисунке~\ref{fig:num}.

\insertfigure{num}{pic/num1-4}{Тестируемые изображенияы}{10cm}

Тест представлен в приложении Г. Для проверки корректности работы нейронной 
сети необходимо сравнить выход результирующего полносвязного слоя и нейронной 
сети. 

На рисунках~\ref{fig:n1} и~\ref{fig:n4} представлены временные диаграммы 
поведенческого моделирования описанной на SystemVerilog нейронной сети и выход 
эталонной нейронной сети на языке Python.

\insertfigurescustom{pic/t1.jpg}{pic/num1_et.png}{16cm}{Временная диаграмма нейронной сети и результат работы эталона на первом изображении}{n1}{stacked}
\insertfigurescustom{pic/t2.jpg}{pic/num4_et.png}{16cm}{Временная диаграмма нейронной сети и результат работы эталона на втором изображении}{n4}{stacked}

В результате проведённого моделирования и сопоставления работы реализованной 
аппаратной модели с эталонной моделью на языке Python с использованием
библиотеки fixpoint было выявлено полное соответствие выходных данных. Это 
подтверждает корректность проектирования и реализации нейронной сети на 
уровне RTL-описания и её работоспособность при выполнении на целевой ПЛИС.

Для оценки ресурсоёмкости и эффективности реализации проекта был проведён 
анализ использования аппаратных ресурсов, таких как логические элементы 
(LUT), триггеры (FF), специализированные арифметические блоки (DSP) и память 
(BRAM), полученные в процессе синтеза HDL-описания. 

Полученные результаты представлены на рисунке~\ref{fig:hw}.

\insertfigure{hw}{pic/hw1.jpg}{Аппаратные затраты}{16cm}

Для дополнительной верификации корректности функционирования разработанной 
системы проводится временная симуляция на основе результатов синтеза и 
имплементации.

На данном этапе осуществляется моделирование с учетом временных задержек, 
рассчитанных инструментом синтеза, однако ещё без учёта конкретных физических 
ресурсов устройства. Постсинтезная симуляция позволяет убедиться, что 
логическая структура проекта сохранена и корректно функционирует после 
преобразования исходного RTL-кода в структуру из примитивов целевой ПЛИС.

Результаты такой симуляции представлены на рисунке~\ref{fig:post-sin}. Видно, 
что формируемые сигналы соответствуют ожидаемым значениям, соблюдаются 
временные зависимости и интерфейсы взаимодействуют корректно.

\insertfigure{post-sin}{pic/p-sin.jpg}{Временная диаграмма нейронной сети на втором изображении при пост-синтез симмуляции}{16cm}

После завершения этапа имплементации, включающего размещение и трассировку 
(place \& route), производится повторная временная симуляция, уже с учетом всех 
физических задержек, связанных с реальной топологией на кристалле. Это 
позволяет оценить окончательные временные характеристики и убедиться, что 
проект не нарушает временных ограничений (timing constraints).

На рисунке~\ref{fig:post-impl} показана временная диаграмма, соответствующая 
данному этапу. Анализ сигналов подтверждает корректную работу всех блоков 
системы, несмотря на добавление реальных задержек межсоединений и элементов 
логики.

\insertfigure{post-impl}{pic/p-impl.jpg}{Временная диаграмма нейронной сети на первом изображении при пост-имплементация симмуляции}{16cm}

В результате этапа имплементации (Placement and Routing) была получена 
финальная топология расположения логических элементов, специализированных 
блоков и соединений на кристалле ПЛИС. Полученная схема представлена на 
рисунке~\ref{fig:impl}.

\insertfigure{impl}{pic/impl.jpg}{Расположение на кристале}{16cm}

На изображении показано физическое распределение ресурсов проекта по кристаллу. 
Ярко выраженные прямоугольные и квадратные блоки соответствуют различным 
логическим компонентам, таким как:

– блоки логики (LUT/FF);

– цифровые сигнальные процессоры (DSP48);

– блоки оперативной памяти (BRAM);

– входные/выходные буферы (IOB).

Размещение было выполнено автоматически средствами среды синтеза и 
имплементации, с учётом топологии кристалла и требований по задержкам сигнала. 
На изображении можно видеть, что распределение элементов достаточно компактное 
и сбалансированное, без чрезмерной фрагментации или перегруженных областей, что 
положительно сказывается на тактовой частоте и стабильности работы системы.

Ещё одним важным параметром при анализе аппаратной реализации является 
потребляемая мощность, которая играет ключевую роль при выборе целевого 
устройства и оценке общей энергоэффективности системы.

На рисунке~\ref{fig:power} представлена оценка энергопотребления 
реализованной нейронной сети на основе результатов имплементации.

\insertfigure{power}{pic/power.jpg}{Оценка потребляемой мощности}{9cm}

Общая мощность составляет 0.166 Вт, из которых:

1 Статическая мощность (Static) — 0.095 Вт, что составляет 57\% от общего 
потребления. Это базовая мощность, потребляемая устройством при включении, вне 
зависимости от его активности.

2 Динамическая мощность (Dynamic) — 0.071 Вт (43\%), обусловлена 
непосредственной работой схем, переключением сигналов, выполнением операций и 
передачей данных.

В структуре динамической мощности:

1 BRAM модули потребляют 0.032 Вт (46\%), что обусловлено активной работой 
памяти при обработке входных данных и весов.

2 DSP блоки — 0.022 Вт (30\%), благодаря интенсивному использованию для 
выполнения операций умножения.

3 Сигналы и тактовые линии составляют 0.008 Вт (11\%) и 0.003 Вт (4\%) 
соответственно.

4 Остальные компоненты (логика и I/O) потребляют сравнительно незначительное 
количество энергии.

Результаты проектирования сгруппированы и представлены в графическом материале 
ГУИР 431289.006 ПЛ.

%------------------------------------------------------------------------------
% \PART 5.5 Text
%------------------------------------------------------------------------------
\section{Устройство осноыных функциональных блоков}
\hspace*{12.5 mm}Ключевым этапом в проектировании аппаратной реализации 
нейросети является синтез и анализ структуры функциональных блоков. Это 
позволяет оценить ресурсоемкость, функциональность и корректность разработанных 
аппаратных модулей, а также выявить узкие места при дальнейшей интеграции в 
систему на базе ПЛИС.

Одним из базовых компонентов системы является блок памяти изображений, 
необходимый для хранения входных данных – изображений рукописных цифр. Данный 
блок реализован с использованием двух логических таблиц, выполняющих формирование 
сигналов WE (Write Enable) и RSTRAM (Reset RAM) для управления чтением и записью. 
Благодаря такой реализации обеспечивается компактность и низкое 
потребление аппаратных ресурсов.

На рисунке~\ref{fig:ram} представлена структурная схема данного модуля. Видно, 
что блок организован эффективно.

\insertfigure{ram}{pic/ram.jpg}{Структура блока памяти изображения}{16cm}

Следующей критически важной частью являются процессорные блоки, отвечающие за 
выполнение вычислений, связанных с прохождением данных через нейронную сеть. 
Эти блоки реализуют операции умножения входных значений на веса, аккумуляции и 
применение функции активации.
Процессорные элементы представлены в двух вариантах.

Первый тип процессорных блоков предназначен для обработки выходного слоя 
нейросети. Его особенностью является наличие встроенной памяти, хранящей 
весовые коэффициенты выходного полносвязного слоя. 

На рисунке~\ref{fig:rc0} показана структура такого блока. В отличие от предыдущего, 
он содержит дополнительные логические цепочки, для чтения значения из третьей
памяти и требует больше логических элиментов.

\insertfigure{rc0}{pic/rco.jpg}{Часть структуры первого блока обработки изображения}{16cm}

На рисунке~\ref{fig:rc} показана часть внутренней структуры второго 
процессорного блока. Видны цепочки логических элементов, осуществляющие 
чтение данных из двух различных блоков памяти весовых коэффициентов и 
выбор одного из коэффициентов для подачи на MAC ядро.

\insertfigure{rc}{pic/rc.jpg}{Часть структуры второго блока обработки изображения}{16cm}

На рисунке~\ref{fig:hw-1} представлены результаты синтеза отдельных данных 
модулей. Каждый блок был синтезирован и проанализирован по количеству 
используемых логических элементов, блоков памяти, триггеров и цифровых 
сигнальных процессоров.

\insertfigurescustom{pic/hw-rco.jpg}{pic/hw-rc.jpg}{8cm}{Аппаратные затраты на различные типы блоков обработки пикселя}{hw-1}{sidebyside}

В частности, использование дополнительной встроенной памяти в блоке второго типа
приводит к увеличению общего количества задействованных BRAM-блоков. Кроме того, 
возросло количество LUT-блоков, что обусловлено добавлением логики формирования
мультиплексоров, обеспечивающих выбор между несколькими блоками хранения весов.

Структура MAC ядра основана на использовании блока DSP48,
что показано на рисунке~\ref{fig:mac}.

\insertfigure{mac}{pic/mac.jpg}{Cтруктурf MAC ядра}{16cm}
